% Options for packages loaded elsewhere
\PassOptionsToPackage{unicode}{hyperref}
\PassOptionsToPackage{hyphens}{url}
\documentclass[
]{article}
\usepackage{xcolor}
\usepackage[margin=1in]{geometry}
\usepackage{amsmath,amssymb}
\setcounter{secnumdepth}{-\maxdimen} % remove section numbering
\usepackage{iftex}
\ifPDFTeX
  \usepackage[T1]{fontenc}
  \usepackage[utf8]{inputenc}
  \usepackage{textcomp} % provide euro and other symbols
\else % if luatex or xetex
  \usepackage{unicode-math} % this also loads fontspec
  \defaultfontfeatures{Scale=MatchLowercase}
  \defaultfontfeatures[\rmfamily]{Ligatures=TeX,Scale=1}
\fi
\usepackage{lmodern}
\ifPDFTeX\else
  % xetex/luatex font selection
\fi
% Use upquote if available, for straight quotes in verbatim environments
\IfFileExists{upquote.sty}{\usepackage{upquote}}{}
\IfFileExists{microtype.sty}{% use microtype if available
  \usepackage[]{microtype}
  \UseMicrotypeSet[protrusion]{basicmath} % disable protrusion for tt fonts
}{}
\makeatletter
\@ifundefined{KOMAClassName}{% if non-KOMA class
  \IfFileExists{parskip.sty}{%
    \usepackage{parskip}
  }{% else
    \setlength{\parindent}{0pt}
    \setlength{\parskip}{6pt plus 2pt minus 1pt}}
}{% if KOMA class
  \KOMAoptions{parskip=half}}
\makeatother
\usepackage{color}
\usepackage{fancyvrb}
\newcommand{\VerbBar}{|}
\newcommand{\VERB}{\Verb[commandchars=\\\{\}]}
\DefineVerbatimEnvironment{Highlighting}{Verbatim}{commandchars=\\\{\}}
% Add ',fontsize=\small' for more characters per line
\usepackage{framed}
\definecolor{shadecolor}{RGB}{248,248,248}
\newenvironment{Shaded}{\begin{snugshade}}{\end{snugshade}}
\newcommand{\AlertTok}[1]{\textcolor[rgb]{0.94,0.16,0.16}{#1}}
\newcommand{\AnnotationTok}[1]{\textcolor[rgb]{0.56,0.35,0.01}{\textbf{\textit{#1}}}}
\newcommand{\AttributeTok}[1]{\textcolor[rgb]{0.13,0.29,0.53}{#1}}
\newcommand{\BaseNTok}[1]{\textcolor[rgb]{0.00,0.00,0.81}{#1}}
\newcommand{\BuiltInTok}[1]{#1}
\newcommand{\CharTok}[1]{\textcolor[rgb]{0.31,0.60,0.02}{#1}}
\newcommand{\CommentTok}[1]{\textcolor[rgb]{0.56,0.35,0.01}{\textit{#1}}}
\newcommand{\CommentVarTok}[1]{\textcolor[rgb]{0.56,0.35,0.01}{\textbf{\textit{#1}}}}
\newcommand{\ConstantTok}[1]{\textcolor[rgb]{0.56,0.35,0.01}{#1}}
\newcommand{\ControlFlowTok}[1]{\textcolor[rgb]{0.13,0.29,0.53}{\textbf{#1}}}
\newcommand{\DataTypeTok}[1]{\textcolor[rgb]{0.13,0.29,0.53}{#1}}
\newcommand{\DecValTok}[1]{\textcolor[rgb]{0.00,0.00,0.81}{#1}}
\newcommand{\DocumentationTok}[1]{\textcolor[rgb]{0.56,0.35,0.01}{\textbf{\textit{#1}}}}
\newcommand{\ErrorTok}[1]{\textcolor[rgb]{0.64,0.00,0.00}{\textbf{#1}}}
\newcommand{\ExtensionTok}[1]{#1}
\newcommand{\FloatTok}[1]{\textcolor[rgb]{0.00,0.00,0.81}{#1}}
\newcommand{\FunctionTok}[1]{\textcolor[rgb]{0.13,0.29,0.53}{\textbf{#1}}}
\newcommand{\ImportTok}[1]{#1}
\newcommand{\InformationTok}[1]{\textcolor[rgb]{0.56,0.35,0.01}{\textbf{\textit{#1}}}}
\newcommand{\KeywordTok}[1]{\textcolor[rgb]{0.13,0.29,0.53}{\textbf{#1}}}
\newcommand{\NormalTok}[1]{#1}
\newcommand{\OperatorTok}[1]{\textcolor[rgb]{0.81,0.36,0.00}{\textbf{#1}}}
\newcommand{\OtherTok}[1]{\textcolor[rgb]{0.56,0.35,0.01}{#1}}
\newcommand{\PreprocessorTok}[1]{\textcolor[rgb]{0.56,0.35,0.01}{\textit{#1}}}
\newcommand{\RegionMarkerTok}[1]{#1}
\newcommand{\SpecialCharTok}[1]{\textcolor[rgb]{0.81,0.36,0.00}{\textbf{#1}}}
\newcommand{\SpecialStringTok}[1]{\textcolor[rgb]{0.31,0.60,0.02}{#1}}
\newcommand{\StringTok}[1]{\textcolor[rgb]{0.31,0.60,0.02}{#1}}
\newcommand{\VariableTok}[1]{\textcolor[rgb]{0.00,0.00,0.00}{#1}}
\newcommand{\VerbatimStringTok}[1]{\textcolor[rgb]{0.31,0.60,0.02}{#1}}
\newcommand{\WarningTok}[1]{\textcolor[rgb]{0.56,0.35,0.01}{\textbf{\textit{#1}}}}
\usepackage{longtable,booktabs,array}
\usepackage{calc} % for calculating minipage widths
% Correct order of tables after \paragraph or \subparagraph
\usepackage{etoolbox}
\makeatletter
\patchcmd\longtable{\par}{\if@noskipsec\mbox{}\fi\par}{}{}
\makeatother
% Allow footnotes in longtable head/foot
\IfFileExists{footnotehyper.sty}{\usepackage{footnotehyper}}{\usepackage{footnote}}
\makesavenoteenv{longtable}
\usepackage{graphicx}
\makeatletter
\newsavebox\pandoc@box
\newcommand*\pandocbounded[1]{% scales image to fit in text height/width
  \sbox\pandoc@box{#1}%
  \Gscale@div\@tempa{\textheight}{\dimexpr\ht\pandoc@box+\dp\pandoc@box\relax}%
  \Gscale@div\@tempb{\linewidth}{\wd\pandoc@box}%
  \ifdim\@tempb\p@<\@tempa\p@\let\@tempa\@tempb\fi% select the smaller of both
  \ifdim\@tempa\p@<\p@\scalebox{\@tempa}{\usebox\pandoc@box}%
  \else\usebox{\pandoc@box}%
  \fi%
}
% Set default figure placement to htbp
\def\fps@figure{htbp}
\makeatother
\setlength{\emergencystretch}{3em} % prevent overfull lines
\providecommand{\tightlist}{%
  \setlength{\itemsep}{0pt}\setlength{\parskip}{0pt}}
\usepackage{bookmark}
\IfFileExists{xurl.sty}{\usepackage{xurl}}{} % add URL line breaks if available
\urlstyle{same}
\hypersetup{
  hidelinks,
  pdfcreator={LaTeX via pandoc}}

\author{}
\date{\vspace{-2.5em}}

\begin{document}

\section{Les réseaux mobiles : une infrastructure
invisible}\label{les-ruxe9seaux-mobiles-une-infrastructure-invisible}

Les réseaux de télécommunications mobiles constituent aujourd'hui l'une
des infrastructures techniques les plus critiques de la société moderne.
Ils permettent la transmission de voix, de données et de messages à
l'échelle mondiale et supportent une part importante des communications
personnelles, industrielles et institutionnelles.

L'architecture actuelle des réseaux mobiles est le résultat d'une
évolution progressive débutée avec les premiers systèmes cellulaires
analogiques, puis avec l'introduction du \textbf{GSM (Global System for
Mobile Communications)} dans les années 1990. Ce standard a profondément
transformé les télécommunications en introduisant une architecture
numérique standardisée, reposant sur une séparation claire entre
plusieurs composants réseau.

Dans un réseau GSM classique, plusieurs éléments principaux coopèrent :

\begin{itemize}
\tightlist
\item
  \textbf{BTS (Base Transceiver Station)} : station radio qui communique
  avec les téléphones mobiles.
\item
  \textbf{BSC (Base Station Controller)} : contrôleur qui gère plusieurs
  BTS et leurs ressources radio.
\item
  \textbf{MSC (Mobile Switching Center)} : commutateur central qui gère
  les appels et la mobilité.
\item
  \textbf{HLR (Home Location Register)} : base de données contenant les
  informations des abonnés.
\item
  \textbf{STP (Signal Transfer Point)} : nœud de routage de la
  signalisation SS7 entre réseaux.
\end{itemize}

Ces composants communiquent entre eux via des protocoles spécialisés,
notamment \textbf{SS7}, un système de signalisation conçu pour
transporter les messages nécessaires au fonctionnement du réseau :
authentification, localisation des abonnés, établissement des appels ou
routage des SMS.

Même si les technologies récentes comme \textbf{4G (LTE)} et \textbf{5G}
ont profondément transformé l'architecture des réseaux mobiles, une
grande partie de cette infrastructure historique reste encore utilisée
aujourd'hui, en particulier dans les couches de compatibilité et dans
certaines applications industrielles ou critiques.

Comprendre ces réseaux nécessite donc d'explorer non seulement leur
fonctionnement radio, mais également l'ensemble de leur infrastructure
de signalisation.

\begin{center}\rule{0.5\linewidth}{0.5pt}\end{center}

\section{Les réseaux mobiles comme systèmes
distribués}\label{les-ruxe9seaux-mobiles-comme-systuxe8mes-distribuuxe9s}

Un réseau mobile est en réalité un \textbf{système distribué complexe}
composé de plusieurs sous-réseaux interconnectés.

Chaque opérateur possède sa propre infrastructure comprenant ses bases
d'abonnés, ses centres de commutation et ses stations radio. Cependant,
les communications doivent également fonctionner entre opérateurs
différents. Pour permettre ces échanges, les réseaux sont reliés via des
mécanismes de signalisation inter-opérateurs.

Dans le monde GSM, cette interconnexion repose principalement sur :

\begin{itemize}
\tightlist
\item
  \textbf{SS7 (Signaling System No.7)}
\item
  \textbf{MAP (Mobile Application Part)}
\item
  \textbf{SIGTRAN / M3UA}, qui permettent de transporter la
  signalisation SS7 sur des réseaux IP modernes.
\end{itemize}

Un élément central dans cette architecture est le \textbf{STP (Signal
Transfer Point)}. Il agit comme un routeur de signalisation capable
d'acheminer les messages entre plusieurs réseaux.

Dans une topologie simplifiée, plusieurs opérateurs peuvent être
connectés à un STP central qui permet d'acheminer les requêtes entre
leurs infrastructures respectives.

Ce type d'architecture permet par exemple :

\begin{itemize}
\tightlist
\item
  le routage d'un \textbf{SMS vers un abonné d'un autre opérateur}
\item
  la localisation d'un téléphone lors d'un \textbf{appel entrant}
\item
  la gestion du \textbf{roaming international}
\end{itemize}

Cette interconnexion rend les réseaux mobiles extrêmement puissants mais
également très complexes à étudier, car elle implique de nombreuses
interactions entre systèmes distincts.

\begin{center}\rule{0.5\linewidth}{0.5pt}\end{center}

\section{Reproduire un réseau télécom en
laboratoire}\label{reproduire-un-ruxe9seau-tuxe9luxe9com-en-laboratoire}

Dans les environnements opérateurs réels, l'infrastructure télécom est
généralement inaccessible aux chercheurs, aux étudiants ou aux
ingénieurs souhaitant comprendre en profondeur le fonctionnement des
réseaux mobiles.

Pour pallier cette difficulté, plusieurs projets open-source ont émergé
afin de reproduire ces architectures en laboratoire. Parmi eux, la suite
\textbf{Osmocom} permet de déployer l'ensemble des composants
nécessaires à la création d'un réseau GSM complet.

Le laboratoire présenté dans cet ouvrage repose sur cette approche. Il
met en place un environnement capable de simuler plusieurs opérateurs
GSM interconnectés via une infrastructure de signalisation complète.

Ce laboratoire inclut notamment :

\begin{itemize}
\tightlist
\item
  plusieurs réseaux GSM indépendants
\item
  un \textbf{STP inter-opérateur}
\item
  la signalisation \textbf{SIGTRAN / M3UA}
\item
  un \textbf{HLR par opérateur}
\item
  un \textbf{MSC}
\item
  un \textbf{BSC}
\item
  une passerelle \textbf{SIP via Asterisk}
\item
  un système de routage \textbf{SMS inter-opérateur}
\end{itemize}

L'objectif de cet environnement est de rendre \textbf{observable et
reproductible} une infrastructure télécom qui est normalement réservée
aux opérateurs.

Dans ce laboratoire, les réseaux peuvent être démarrés, modifiés et
analysés en temps réel. Les échanges de signalisation peuvent être
inspectés, les routes modifiées et les interactions entre opérateurs
étudiées.

Par exemple, il est possible d'observer comment une requête de
localisation d'abonné traverse plusieurs réseaux avant d'atteindre le
HLR correspondant, ou comment un SMS est routé d'un opérateur vers un
autre via un STP central.

\begin{center}\rule{0.5\linewidth}{0.5pt}\end{center}

\subsection{Architecture simplifiée du
laboratoire}\label{architecture-simplifiuxe9e-du-laboratoire}

Chaque opérateur possède ses propres identifiants réseau et son propre
plan de signalisation. Le STP central agit comme un point d'échange
permettant d'acheminer les messages entre ces différents réseaux.

Ce type d'environnement permet de reproduire des scénarios réalistes
tout en conservant un contrôle complet sur l'infrastructure.

\begin{center}\includegraphics[width=9.47in]{../../../../tmp/RtmpZY6wjN/file26cc44d9761d9} \end{center}

\begin{center}\rule{0.5\linewidth}{0.5pt}\end{center}

\section{Implémentation du réseau SS7 et analyse des
dumps}\label{impluxe9mentation-du-ruxe9seau-ss7-et-analyse-des-dumps}

Dans un réseau télécom réel, la signalisation est le système nerveux de
l'infrastructure. Les appels, les SMS et la gestion des abonnés ne
reposent pas uniquement sur la radio : ils dépendent d'un réseau de
signalisation distinct chargé de coordonner les équipements du cœur de
réseau.

Dans l'architecture GSM classique, ce rôle est assuré par \textbf{SS7
(Signalling System No.~7)}. Historiquement conçu pour les réseaux
téléphoniques commutés, SS7 a ensuite été adapté aux réseaux mobiles
pour transporter des protocoles comme \textbf{MAP}, utilisés pour la
gestion des abonnés et des services mobiles.

Dans ce laboratoire, la signalisation SS7 n'est pas simulée de manière
abstraite : elle est réellement implémentée sur IP à l'aide de
\textbf{SIGTRAN et M3UA}, ce qui permet d'observer et d'analyser les
échanges réels entre les composants du réseau. Le projet reproduit ainsi
une infrastructure complète d'interconnexion entre plusieurs réseaux
mobiles simulés.

\begin{center}\rule{0.5\linewidth}{0.5pt}\end{center}

\subsection{Le rôle du réseau SS7 dans l'architecture
GSM}\label{le-ruxf4le-du-ruxe9seau-ss7-dans-larchitecture-gsm}

Dans un réseau GSM, la radio (BTS et BSC) ne constitue qu'une partie de
l'infrastructure. Les décisions concernant l'acheminement des appels ou
des SMS sont prises dans le cœur de réseau.

Les principaux éléments impliqués sont :

\begin{itemize}
\tightlist
\item
  \textbf{MSC (Mobile Switching Center)} : Centre de commutation qui
  gère les appels et la mobilité.
\item
  \textbf{HLR (Home Location Register)} : Base de données contenant les
  informations des abonnés.
\item
  \textbf{STP (Signal Transfer Point)} : Routeur de signalisation SS7
  chargé d'acheminer les messages.
\end{itemize}

Dans une architecture classique, tous ces éléments échangent des
messages SS7 pour coordonner leurs actions.

Un SMS, par exemple, implique généralement plusieurs échanges :

\begin{enumerate}
\def\labelenumi{\arabic{enumi}.}
\tightlist
\item
  le MSC reçoit le message du téléphone
\item
  il interroge le HLR pour trouver le destinataire
\item
  le message est transmis au SMSC
\item
  le SMSC demande au HLR où se trouve l'abonné
\item
  le message est finalement délivré au MSC du réseau cible
\end{enumerate}

Chaque étape correspond à une série de messages MAP transportés sur SS7.

\begin{center}\rule{0.5\linewidth}{0.5pt}\end{center}

\subsection{SS7 sur IP : SIGTRAN et
M3UA}\label{ss7-sur-ip-sigtran-et-m3ua}

Les réseaux modernes n'utilisent plus les liaisons SS7 traditionnelles
sur TDM. La signalisation est aujourd'hui transportée sur IP grâce à une
famille de protocoles appelée \textbf{SIGTRAN}.

Dans cette architecture, le transport se fait généralement via :

\begin{itemize}
\tightlist
\item
  \textbf{SCTP} pour la couche transport
\item
  \textbf{M3UA (MTP3 User Adaptation)} pour transporter les messages SS7
  sur IP
\end{itemize}

Dans le laboratoire, chaque composant Osmocom (MSC, BSC, STP) communique
via \textbf{M3UA sur SCTP}, exactement comme dans les déploiements
opérateurs modernes.

Cela permet d'observer directement les flux de signalisation dans un
analyseur réseau comme Wireshark.

\begin{center}\includegraphics[width=2.31in]{../../../../tmp/RtmpZY6wjN/file26cc4743cb264} \end{center}

\begin{center}\rule{0.5\linewidth}{0.5pt}\end{center}

\subsection{L'architecture
multi-opérateurs}\label{larchitecture-multi-opuxe9rateurs}

L'une des particularités du projet est de simuler plusieurs réseaux
mobiles indépendants, chacun possédant son propre cœur de réseau.

Chaque opérateur est isolé dans un conteneur Docker contenant :

\begin{itemize}
\tightlist
\item
  un \textbf{MSC}
\item
  un \textbf{HLR}
\item
  un \textbf{BSC}
\item
  un \textbf{STP local}
\item
  un \textbf{MGW}
\item
  un \textbf{SMSC}
\end{itemize}

Tous ces composants communiquent localement via l'interface
\textbf{127.0.0.1}, ce qui évite les problèmes de synchronisation au
démarrage des conteneurs.

L'interconnexion entre les différents opérateurs est assurée par un
composant central : l'\textbf{Inter-STP}.

Ce nœud agit comme un \textbf{routeur SS7 inter-opérateurs}, chargé
d'acheminer la signalisation entre les différents réseaux.

\begin{center}\includegraphics[width=5.17in]{../../../../tmp/RtmpZY6wjN/file26cc471e4035a} \end{center}

\begin{center}\rule{0.5\linewidth}{0.5pt}\end{center}

\subsection{Topologie réseau Docker}\label{topologie-ruxe9seau-docker}

La topologie réseau du laboratoire repose sur deux types de réseaux
Docker distincts :

\begin{enumerate}
\def\labelenumi{\arabic{enumi}.}
\tightlist
\item
  un \textbf{réseau backbone inter-opérateurs}, utilisé pour transporter
  la signalisation SS7
\item
  un \textbf{réseau privé par opérateur}, utilisé pour le trafic interne
  GSM
\end{enumerate}

Cette séparation reflète l'organisation réelle des réseaux télécom où la
signalisation circule sur un réseau distinct du trafic utilisateur.

\subsubsection{Réseaux}\label{ruxe9seaux}

\begin{longtable}[]{@{}
  >{\raggedright\arraybackslash}p{(\linewidth - 4\tabcolsep) * \real{0.1940}}
  >{\raggedright\arraybackslash}p{(\linewidth - 4\tabcolsep) * \real{0.2239}}
  >{\raggedright\arraybackslash}p{(\linewidth - 4\tabcolsep) * \real{0.5821}}@{}}
\toprule\noalign{}
\begin{minipage}[b]{\linewidth}\raggedright
Réseau Docker
\end{minipage} & \begin{minipage}[b]{\linewidth}\raggedright
Plage
\end{minipage} & \begin{minipage}[b]{\linewidth}\raggedright
Rôle
\end{minipage} \\
\midrule\noalign{}
\endhead
\bottomrule\noalign{}
\endlastfoot
\texttt{gsm-inter} & \texttt{172.20.0.0/24} & Backbone interop (M3UA
inter-STP) \\
\texttt{gsm-net-opN} & \texttt{172.20.N.0/24} & Réseau privé opérateur N
(GSMTAP, GPRS) \\
\end{longtable}

Le réseau \textbf{gsm-inter} joue le rôle de backbone de signalisation.
C'est sur ce réseau que transitent les connexions \textbf{SCTP / M3UA}
entre les STP opérateurs et l'Inter-STP.

Chaque opérateur possède en parallèle un réseau privé dédié.

\subsubsection{Attribution des adresses
IP}\label{attribution-des-adresses-ip}

Pour simplifier la génération automatique des configurations,
l'adressage IP suit une formule simple.

\begin{longtable}[]{@{}
  >{\raggedright\arraybackslash}p{(\linewidth - 4\tabcolsep) * \real{0.2063}}
  >{\raggedright\arraybackslash}p{(\linewidth - 4\tabcolsep) * \real{0.3968}}
  >{\raggedright\arraybackslash}p{(\linewidth - 4\tabcolsep) * \real{0.3968}}@{}}
\toprule\noalign{}
\begin{minipage}[b]{\linewidth}\raggedright
Composant
\end{minipage} & \begin{minipage}[b]{\linewidth}\raggedright
IP backbone (\texttt{gsm-inter})
\end{minipage} & \begin{minipage}[b]{\linewidth}\raggedright
IP privée (\texttt{gsm-net-opN})
\end{minipage} \\
\midrule\noalign{}
\endhead
\bottomrule\noalign{}
\endlastfoot
Inter-STP & \texttt{172.20.0.10} & --- \\
Conteneur OpN & \texttt{172.20.0.(10+N)} & \texttt{172.20.N.10} \\
\end{longtable}

La règle générale est la suivante :

\begin{verbatim}
opérateur N
backbone : 172.20.0.(10+N)
réseau privé : 172.20.N.10
\end{verbatim}

Cette convention permet de générer automatiquement toutes les
configurations réseau lors du démarrage du laboratoire.

\subsubsection{Communication
intra-conteneur}\label{communication-intra-conteneur}

À l'intérieur d'un conteneur opérateur, tous les composants Osmocom
communiquent via \textbf{localhost}.

Ce choix présente deux avantages :

\begin{itemize}
\tightlist
\item
  éliminer les \textbf{conditions de course Docker} lors du démarrage
\item
  garantir que le STP est disponible immédiatement pour les autres
  composants
\end{itemize}

Le STP local agit ainsi comme \textbf{hub de signalisation
intra-opérateur}.

\subsubsection{Communication
inter-conteneurs}\label{communication-inter-conteneurs}

La communication entre opérateurs est volontairement limitée.

Le seul lien qui traverse réellement le réseau Docker est la connexion
entre :

\begin{itemize}
\tightlist
\item
  le \textbf{STP opérateur}
\item
  l'\textbf{Inter-STP}
\end{itemize}

\begin{center}\includegraphics[width=5.47in]{../../../../tmp/RtmpZY6wjN/file26cc44ce2effe} \end{center}

\begin{center}\rule{0.5\linewidth}{0.5pt}\end{center}

\section{Signalisation SS7
détaillée}\label{signalisation-ss7-duxe9tailluxe9e}

\subsection{Point codes}\label{point-codes}

Dans SS7, chaque nœud possède un identifiant unique appelé \textbf{point
code}.

Ces identifiants permettent de déterminer vers quel équipement doit être
routé un message.

Le laboratoire utilise le format \textbf{ITU 14-bit} :

\begin{verbatim}
zone.network.node
\end{verbatim}

\begin{longtable}[]{@{}lll@{}}
\toprule\noalign{}
Nœud & Point Code & Formule \\
\midrule\noalign{}
\endhead
\bottomrule\noalign{}
\endlastfoot
Inter-STP & \texttt{0.23.0} & fixe \\
MSC OpN & \texttt{N.23.1} & zone=N \\
STP OpN & \texttt{N.23.2} & zone=N \\
BSC OpN & \texttt{N.23.3} & zone=N \\
\end{longtable}

Cette convention permet d'identifier immédiatement l'opérateur et le
type d'équipement.

\subsection{Routing Contexts (RCTX)}\label{routing-contexts-rctx}

Dans M3UA, le \textbf{Routing Context} permet d'identifier un flux de
signalisation.

Le projet utilise une formule systématique pour générer ces identifiants
:

\begin{longtable}[]{@{}lll@{}}
\toprule\noalign{}
RCTX & Formule & Usage \\
\midrule\noalign{}
\endhead
\bottomrule\noalign{}
\endlastfoot
\texttt{N×100+10} & \texttt{rctx\_msc} & MSC OpN → STP OpN \\
\texttt{N×100+20} & \texttt{rctx\_stp} & STP OpN \\
\texttt{N×100+30} & \texttt{rctx\_bsc} & BSC OpN \\
\textbf{\texttt{N×100+50}} & \textbf{\texttt{rctx\_inter}} & \textbf{STP
OpN → Inter-STP} \\
\end{longtable}

Le contexte le plus important est \textbf{rctx\_inter}, utilisé pour
l'interconnexion entre opérateurs.

Une incohérence dans ce contexte empêche la signalisation inter-réseau
de fonctionner correctement.

\subsection{Configuration de
l'Inter-STP}\label{configuration-de-linter-stp}

L'Inter-STP est le routeur central de signalisation.

Il utilise une configuration simplifiée : chaque opérateur est
représenté par un \textbf{Application Server (AS)}.

Le routage repose uniquement sur le \textbf{Destination Point Code}.

\begin{center}\rule{0.5\linewidth}{0.5pt}\end{center}

\subsection{Parcours complet d'un message
SS7}\label{parcours-complet-dun-message-ss7}

Le chemin d'un message de signalisation peut être représenté de la
manière suivante :

\begin{center}\includegraphics[width=6.86in]{../../../../tmp/RtmpZY6wjN/file26cc41285f9e1} \end{center}

Le message suit donc quatre étapes :

\begin{enumerate}
\def\labelenumi{\arabic{enumi}.}
\tightlist
\item
  le \textbf{MSC} envoie le message au \textbf{STP local}
\item
  le STP local transmet le message à l'\textbf{Inter-STP}
\item
  l'Inter-STP sélectionne l'opérateur cible
\item
  le STP distant délivre le message au \textbf{MSC destinataire}
\end{enumerate}

Ce mécanisme reproduit exactement la logique d'interconnexion des
réseaux mobiles.

\begin{center}\rule{0.5\linewidth}{0.5pt}\end{center}

\section{Observabilité et
diagnostic}\label{observabilituxe9-et-diagnostic}

Une fois l'infrastructure SS7 en place, l'un des aspects les plus
intéressants du laboratoire est la possibilité d'observer la
signalisation en temps réel.

Dans les réseaux opérateurs, l'accès à la signalisation SS7 est
généralement extrêmement restreint. Les équipements de cœur de réseau
sont isolés et les ingénieurs disposent seulement d'outils internes pour
analyser les flux.

Dans ce laboratoire, la situation est radicalement différente : toute la
signalisation peut être inspectée.

Deux mécanismes principaux permettent cette observation :

\begin{itemize}
\tightlist
\item
  les \textbf{interfaces VTY des composants Osmocom}
\item
  les \textbf{captures réseau avec Wireshark}
\end{itemize}

Les interfaces VTY exposent l'état interne des nœuds SS7. Elles
permettent par exemple de vérifier les associations M3UA actives, la
table de routage ou l'état des routes de signalisation.

Une commande simple permet par exemple de consulter les connexions ASP :

\begin{verbatim}
show cs7 instance 0 asp
\end{verbatim}

Cette commande indique si les associations de signalisation sont
\textbf{actives}, \textbf{inactives} ou \textbf{en attente}.

Une autre commande essentielle est :

\begin{verbatim}
show cs7 instance 0 route
\end{verbatim}

Elle affiche la table de routage SS7 et permet de vérifier si les routes
sont disponibles.

Dans un système sain, les routes dynamiques correspondant aux MSC et aux
BSC apparaissent comme \textbf{available}. Si une route apparaît comme
\textbf{PROHIB}, cela signifie que la signalisation ne peut pas
atteindre ce point code.

Ces informations constituent ce que l'on appelle des \textbf{dumps de
signalisation}. Elles permettent de comprendre précisément comment le
réseau prend ses décisions de routage.

\subsection{Analyse des flux avec
Wireshark}\label{analyse-des-flux-avec-wireshark}

Le laboratoire permet également de capturer directement les flux de
signalisation transportés sur le réseau Docker.

Les protocoles visibles sont les suivants :

\begin{itemize}
\tightlist
\item
  \textbf{SCTP} --- transport utilisé par SIGTRAN
\item
  \textbf{M3UA} --- encapsulation des messages SS7
\item
  \textbf{SCCP} --- routage de signalisation
\item
  \textbf{MAP} --- gestion des services mobiles
\end{itemize}

Dans Wireshark, quelques filtres simples permettent d'isoler ces
protocoles :

\begin{verbatim}
m3ua
sccp
gsm_map
sctp.srcport == 2908
\end{verbatim}

Ces captures permettent d'observer des procédures fondamentales du
réseau GSM :

\begin{itemize}
\tightlist
\item
  requêtes de localisation d'abonné
\item
  routage des SMS
\item
  messages de gestion de mobilité
\item
  interrogations du HLR
\end{itemize}

Le laboratoire devient ainsi un véritable \textbf{microscope de
signalisation télécom}.

\begin{center}\rule{0.5\linewidth}{0.5pt}\end{center}

\section{Diagrammes UML de
l'architecture}\label{diagrammes-uml-de-larchitecture}

\subsection{1. Diagramme de classes : Composants
SS7}\label{diagramme-de-classes-composants-ss7}

\begin{center}\includegraphics[width=5.72in]{../../../../tmp/RtmpZY6wjN/file26cc44c0c5b24} \end{center}

\begin{center}\rule{0.5\linewidth}{0.5pt}\end{center}

\subsection{2. Flux d'appel
inter-opérateurs}\label{flux-dappel-inter-opuxe9rateurs}

\begin{center}\includegraphics[width=2.27in]{../../../../tmp/RtmpZY6wjN/file26cc4632540c6} \end{center}

\begin{center}\rule{0.5\linewidth}{0.5pt}\end{center}

\subsection{3. Stack de protocoles SS7}\label{stack-de-protocoles-ss7}

\begin{center}\includegraphics[width=3.4in]{../../../../tmp/RtmpZY6wjN/file26cc43e96132e} \end{center}

\begin{center}\rule{0.5\linewidth}{0.5pt}\end{center}

\subsection{4. SMS inter-opérateurs}\label{sms-inter-opuxe9rateurs}

\begin{center}\includegraphics[width=8.37in]{../../../../tmp/RtmpZY6wjN/file26cc42317974d} \end{center}

\begin{center}\rule{0.5\linewidth}{0.5pt}\end{center}

\section{Conclusion}\label{conclusion}

Les réseaux mobiles modernes reposent sur une architecture complexe dans
laquelle la radio n'est qu'une partie visible de l'infrastructure.

Derrière chaque appel ou chaque SMS se trouve un réseau de signalisation
chargé de coordonner les équipements du cœur de réseau. SS7 a longtemps
constitué l'épine dorsale de cette infrastructure, et bien que les
réseaux évoluent vers des architectures IP plus récentes, les principes
fondamentaux de la signalisation restent similaires.

Le projet \textbf{osmo\_egprs} propose une approche différente de
l'apprentissage des télécommunications : au lieu d'étudier les
protocoles de manière abstraite, il permet de \textbf{reconstruire un
réseau complet}, puis d'observer son fonctionnement interne.

Grâce à la conteneurisation et à l'automatisation de la configuration,
il devient possible de déployer en quelques minutes un environnement qui
reproduit :

\begin{itemize}
\tightlist
\item
  plusieurs réseaux mobiles indépendants
\item
  un backbone de signalisation SS7
\item
  des échanges inter-opérateurs
\item
  des services voix et SMS
\end{itemize}

Ce laboratoire transforme ainsi un système habituellement inaccessible
en un environnement expérimental ouvert.

Il permet non seulement de comprendre l'architecture des réseaux
mobiles, mais aussi d'explorer leur fonctionnement réel --- un pas
essentiel pour quiconque souhaite étudier les télécommunications, leur
histoire et leur sécurité.

\begin{center}\rule{0.5\linewidth}{0.5pt}\end{center}

\section{Annexe : SDR \& Matériel (UHD,
BladeRF)}\label{annexe-sdr-matuxe9riel-uhd-bladerf}

\subsection{UHD (USRP/Ettus)}\label{uhd-usrpettus}

\begin{itemize}
\item
  \textbf{libuhd-dev, uhd-host} : Drivers et outils pour USRP.
\item
  \textbf{Images firmware} :
\end{itemize}

\begin{Shaded}
\begin{Highlighting}[]
\ExtensionTok{python3}\NormalTok{ /usr/lib/uhd/utils/uhd\_images\_downloader.py}
\end{Highlighting}
\end{Shaded}

\begin{itemize}
\tightlist
\item
  \textbf{Règle udev} pour détection auto :
\end{itemize}

\begin{Shaded}
\begin{Highlighting}[]
\FunctionTok{wget}\NormalTok{ https://raw.githubusercontent.com/EttusResearch/uhd/refs/heads/master/host/utils/uhd{-}usrp.rules}
\FunctionTok{cp}\NormalTok{ uhd{-}usrp.rules /etc/udev/rules.d/}
\end{Highlighting}
\end{Shaded}

\subsection{BladeRF (Nuand)}\label{bladerf-nuand}

\begin{itemize}
\tightlist
\item
  \textbf{libbladerf-dev} + build depuis les sources :
\end{itemize}

\begin{Shaded}
\begin{Highlighting}[]
\FunctionTok{git}\NormalTok{ clone https://github.com/nuand/BladeRF /opt/GSM/BladeRF}
\BuiltInTok{cd}\NormalTok{ /opt/GSM/BladeRF }\KeywordTok{\&\&} \FunctionTok{mkdir}\NormalTok{ build }\KeywordTok{\&\&} \BuiltInTok{cd}\NormalTok{ build}
\FunctionTok{cmake}\NormalTok{ .. }\KeywordTok{\&\&} \FunctionTok{make} \AttributeTok{{-}j}\VariableTok{$(}\FunctionTok{nproc}\VariableTok{)} \KeywordTok{\&\&} \FunctionTok{make}\NormalTok{ install }\KeywordTok{\&\&} \ExtensionTok{ldconfig}
\end{Highlighting}
\end{Shaded}

\begin{center}\rule{0.5\linewidth}{0.5pt}\end{center}

\section{Librairies Osmocom de base}\label{librairies-osmocom-de-base}

Ces librairies sont la fondation de toute la stack GSM Osmocom.

\begin{itemize}
\tightlist
\item
  \textbf{libosmocore} : Primitives GSM, timers, gestion protocoles.
\item
  \textbf{libosmo-netif} : Interfaces réseau.
\item
  \textbf{libosmo-abis} : Support interface Abis (BTS ↔ BSC).
\end{itemize}

\textbf{Build de chaque librairie :}

\begin{Shaded}
\begin{Highlighting}[]
\FunctionTok{git}\NormalTok{ clone https://github.com/osmocom/}\OperatorTok{\textless{}}\NormalTok{repo}\OperatorTok{\textgreater{}}
\BuiltInTok{cd} \OperatorTok{\textless{}}\NormalTok{repo}\OperatorTok{\textgreater{}} \KeywordTok{\&\&} \ExtensionTok{autoreconf} \AttributeTok{{-}fi} \KeywordTok{\&\&} \ExtensionTok{./configure} \KeywordTok{\&\&} \FunctionTok{make} \AttributeTok{{-}j}\VariableTok{$(}\FunctionTok{nproc}\VariableTok{)} \KeywordTok{\&\&} \FunctionTok{make}\NormalTok{ install }\KeywordTok{\&\&} \ExtensionTok{ldconfig}
\end{Highlighting}
\end{Shaded}

\emph{(remplacer \texttt{\textless{}repo\textgreater{}} par le nom de la
librairie)}

\begin{center}\rule{0.5\linewidth}{0.5pt}\end{center}

\emph{Document généré le 2026-03-06 avec R Markdown et Graphviz}

\end{document}
